\chapter{Реализация и анализ эффективности программного комплекса}
\label{ch:implementation}

\section{Реализация модулей мониторинга, тестирования и интерфейса}

Реализация комплекса построена вокруг идеи «одного окна» для системного анализа, сформулированной на этапе проектирования архитектуры. На практике это означает, что пользователь должен иметь возможность одновременно наблюдать текущее состояние системы, запускать измерительные сценарии и интерпретировать результаты в одном и том же приложении, не переходя между несколькими разрозненными инструментами. С инженерной точки зрения такая постановка задачи существенно усложняет реализацию, поскольку мониторинг и тестирование конкурируют за вычислительные ресурсы, а сам пользовательский интерфейс при неудачной организации может стать дополнительным источником искажения измерений.

Если вычислительные тесты выполняются в том же потоке, что и отрисовка графического интерфейса, приложение быстро теряет отзывчивость, пользователь лишается обратной связи о ходе измерения, а мониторинг перестаёт быть полезным контекстом. Если же телеметрия обновляется слишком часто или избыточно тяжёлым способом, она начинает оказывать заметное влияние на результаты бенчмарков, особенно на короткие тесты, чувствительные к системным возмущениям. Следовательно, реализация комплекса должна была решить две взаимосвязанные задачи: изолировать измерительную логику от слоя представления и одновременно обеспечить их координацию в рамках единой интерактивной среды.

В результате программный комплекс был реализован как совокупность модулей, взаимодействующих через понятные и достаточно узкие интерфейсы. Подсистема мониторинга отвечает за периодический сбор и предварительную нормализацию метрик, подсистема бенчмарков --- за выполнение тестов и статистическую агрегацию результатов, а графический слой --- за запуск сценариев, отображение текущего состояния и интерпретацию итоговых значений. Такая схема обеспечивает не только функциональную полноту приложения, но и его сопровождаемость: логика получения данных, вычислительная нагрузка и пользовательское представление не смешиваются в одном уровне абстракции.

\subsection{Подсистема мониторинга: сбор и нормализация метрик}

\textbf{Сбор телеметрии.} Модуль \texttt{pkg/profiling} реализует набор функций получения метрик преимущественно на базе библиотеки \texttt{gopsutil}, которая предоставляет переносимые интерфейсы для доступа к системной информации. Для центрального процессора выполняется получение сведений о модели, частотных характеристиках и текущей загрузке каждого логического ядра. Последняя измеряется с использованием \texttt{cpu.Percent} на интервале порядка одной секунды. Выбор именно такого окна наблюдения обусловлен компромиссом между оперативностью и устойчивостью: слишком короткий интервал давал бы излишне шумные значения, тогда как слишком длинный ухудшал бы восприятие динамики нагрузки в интерфейсе.

Для подсистемы памяти используется \texttt{mem.VirtualMemory}, что позволяет получить сведения о доступной, используемой и общей памяти, а также о ряде вспомогательных характеристик, связанных с состоянием памяти на конкретной платформе. Для дисковой подсистемы применяется агрегирование счётчиков \texttt{disk.IOCounters} и вычисление скоростей чтения и записи как разности между последовательными измерениями. Аналогичный принцип используется для сети через \texttt{net.IOCounters}, где интерес представляют не абсолютные накопленные счётчики, а текущая интенсивность обмена. Такой подход переводит технические системные величины в форму, более понятную пользователю и пригодную для интерпретации в контексте прикладной нагрузки.

Отдельно реализован сбор общей системной информации: времени работы системы, версии платформы, а также средней нагрузки (\texttt{load average}) там, где она доступна. Показатель средней нагрузки дополняет проценты загрузки CPU и позволяет оценивать не только факт использования процессорного времени, но и степень насыщения планировщика. Если значения нагрузки стабильно превышают число логических процессоров, это может свидетельствовать о наличии очереди готовых к выполнению задач и о дефиците вычислительных ресурсов при текущем профиле работы.

Важной практической частью мониторинга является показ «топ-процессов». В модуле \texttt{pkg/profiling/processes.go} формируется список процессов, для которых запрашиваются показатели использования CPU и памяти. Затем результаты сортируются и представляются пользователю в текстовой таблице фиксированной ширины. Несмотря на внешнюю простоту, данная функция играет существенную роль в интерпретации измерений. Она позволяет не просто увидеть, что система нагружена, а понять, какими именно процессами формируется эта нагрузка. Например, нестабильный результат теста может быть объяснён работой компилятора, индексацией IDE, фоновым обновлением ОС или активностью браузера.

С инженерной точки зрения подсистема мониторинга решает не только задачу получения метрик, но и задачу их нормализации. Пользователю важны не внутренние системные счётчики сами по себе, а прикладные показатели в естественных единицах измерения. Поэтому модуль преобразует агрегированные значения в формы, пригодные для непосредственного восприятия: для диска и сети --- в скорости обмена, для CPU --- в проценты загрузки по ядрам, для памяти --- в объёмные характеристики с интерпретируемым отношением занятости и доступности ресурсов.

Кроме того, при реализации мониторинга учитывалась необходимость сохранения умеренной собственной нагрузки инструмента. Если телеметрия будет собираться чрезмерно часто или с чрезмерной детализацией, комплекс начнёт влиять на систему не только как наблюдатель, но и как активный потребитель ресурсов. Поэтому период обновления и состав отображаемых данных были выбраны так, чтобы обеспечить полезную информативность без заметного ухудшения чистоты измерений.

\subsection{Подсистема бенчмарков: контракт, каталог и механизм запуска}

\textbf{Организация запуска бенчмарков.} Подсистема тестирования построена на принципе унифицированного контракта. Все тесты приводятся к единому интерфейсу \texttt{BenchmarkFunc}, возвращающему числовой показатель и строку единицы измерения. Такое решение имеет фундаментальное значение для всей системы, так как позволяет абстрагировать пользовательский интерфейс от внутренней специфики конкретного алгоритма. Панель отображения результатов не обязана знать, что именно выполнялось внутри теста: сортировка, шифрование, вычисление хеша или операция с памятью. Ей достаточно получить итоговое значение и единицу измерения, после чего результат может быть представлен в общей форме и дополнен интерпретацией.

В графическом слое каждая задача описывается структурой \texttt{BenchmarkTask}, включающей название, категорию, краткое и научное описание, функцию интерпретации результата и ссылку на сам тест. Благодаря этому набор бенчмарков превращается в формализованный каталог, который можно расширять без изменения базового механизма отображения. С архитектурной точки зрения это позволяет отделить «что измеряется» от «как показывается», а с практической --- упрощает развитие комплекса и снижает стоимость добавления новых сценариев.

В рамках реализации особое значение имеет разделение тестов по категориям. Пользователь воспринимает измерения не как набор низкоуровневых функций, а как группы задач: «Процессор», «Память», «Математика», «Криптография». Такой подход решает сразу две задачи. Во-первых, он уменьшает когнитивную нагрузку, позволяя запускать измерения в понятных прикладных терминах. Во-вторых, он упрощает последующий анализ, поскольку результаты уже сгруппированы по профилю нагрузки и могут быть сопоставлены внутри своей функциональной области.

С инженерной точки зрения каталог бенчмарков представляет собой промежуточный слой между вычислительной логикой и интерфейсом. Он позволяет описывать тест не только как функцию, но и как пользовательский сценарий с определённым смыслом. Благодаря этому система приобретает свойства, характерные не просто для набора утилит, а для целостного аналитического инструмента.

\subsection{Статистическая обработка результатов и устойчивость к шуму}

Особое внимание в реализации уделено статистике, так как измерения выполняются в условиях многозадачной операционной системы и неизбежно подвержены шуму. Даже если тестовая функция детерминирована, фактическое время её выполнения зависит от множества внешних факторов: миграции потоков между ядрами, фоновых прерываний, активности других процессов, изменений частоты процессора, прогрева кэшей и иных воздействий. По этой причине единичный запуск не может рассматриваться как достаточная основа для надёжного вывода.

\texttt{TestRunner} использует следующую схему:
\begin{enumerate}
    \item выполняется разогрев, чтобы исключить влияние ленивой инициализации и прогрева кэшей;
    \item выполняется серия итераций (минимум 10), результаты сортируются;
    \item применяется IQR-фильтрация для удаления выбросов, возникающих из-за планировщика, фоновых прерываний и коротких всплесков активности других процессов;
    \item рассчитываются робастные показатели: медиана, MAD и коэффициент вариации, после чего итоговое значение формируется как взвешенная комбинация медианы и среднего.
\end{enumerate}

Разогрев необходим для того, чтобы исключить влияние неустоявшегося начального состояния. В первых итерациях часть времени может тратиться на выделение памяти, инициализацию структур данных, наполнение кэшей и стабилизацию режима исполнения. Если эти значения включить в итог без поправок, они могут исказить реальную оценку устойчивой производительности. Серия повторных измерений, в свою очередь, даёт возможность оценить разброс и выявить нестабильность, связанную с внешними факторами.

Использование IQR-фильтрации позволяет отбрасывать редкие выбросы, возникающие, например, при внезапной активности планировщика или кратковременном вмешательстве фоновой задачи. Такой подход особенно полезен в интерактивном приложении, где важно получать практично устойчивые результаты без необходимости переходить к сложным статистическим моделям, избыточным для повседневного системного анализа.

Коэффициент вариации, выраженный в процентах, выводится пользователю как индикатор качества измерения. Чем выше разброс, тем менее надёжны единичные значения и тем сильнее влияние внешних факторов. Такая интерпретация особенно важна при сравнении устройств, поиске регресса производительности и оценке воспроизводимости эксперимента. Пользователь получает не только итоговый показатель, но и информацию о том, насколько этому показателю можно доверять в текущих условиях.

Выбранная методика удобна тем, что сохраняет баланс между практичностью и строгостью. С одной стороны, измерения выполняются достаточно быстро для интерактивного применения и не превращают запуск теста в длительную лабораторную процедуру. С другой стороны, пользователь получает не просто «одно число», а набор статистически осмысленных характеристик, достаточных для принятия решения о повторном прогоне, изменении условий эксперимента или сравнении результатов между платформами.

\subsection{Графический интерфейс и конкуррентность}

\textbf{Графический интерфейс.} Точка входа приложения (\texttt{cmd/analyzer/main.go}) создаёт окно и набор вкладок: «Мониторинг», «Процессор», «Память», «Математика», «Криптография». Такая организация пространства интерфейса отражает принятую на этапе проектирования классификацию сценариев анализа. Пользователь получает не перегруженное единое полотно, а структурированную среду, в которой каждой группе задач соответствует собственная логическая область.

Вкладка мониторинга построена на панели \texttt{NewDashboardPanel}, внутри которой в фоне выполняется цикл с периодом около одной секунды, собирающий метрики и обновляющий виджеты. Визуализация загрузки по ядрам реализована через набор индикаторов \texttt{ProgressBar}, создаваемых динамически под фактическое число логических процессоров. Такой подход позволяет адаптировать интерфейс к разным платформам без жёстко зафиксированного числа ядер и одновременно даёт пользователю наглядное представление о распределении нагрузки, а не только о её усреднённом значении.

Для запуска измерений используются панели \texttt{NewBenchmarkPanel(category)}, которые инкапсулируют логику работы с каталогом тестов в конкретной категории. Панель отвечает за запуск серии измерений, ожидание результата, отображение прогресса и выдачу итогового значения вместе с интерпретацией. Подобное решение обеспечивает единообразие пользовательского опыта: независимо от категории нагрузок пользователь взаимодействует с логически одинаковыми элементами интерфейса.

Отдельным элементом является интерактивная криптопанель \texttt{NewInteractiveCryptoPanel}. Она позволяет выбирать алгоритм, параметры ключа или хеша, а также объём тестовых данных. Включение такой панели в состав комплекса имеет особую ценность, поскольку переводит приложение из режима простого «запуска предопределённых тестов» в режим исследовательского взаимодействия с вычислительной системой. Пользователь может варьировать параметры и непосредственно наблюдать, как эти изменения отражаются на временных затратах и пропускной способности.

Для больших объёмов данных применяется сэмплирование и экстраполяция времени, что делает интерфейс практичным и предотвращает длительную блокировку пользователя при попытке выполнить крайне тяжёлую операцию в тестовом режиме. С инженерной точки зрения это пример компромисса между точностью и удобством: приложение сохраняет возможность оценивать масштаб поведения алгоритма без необходимости реально выполнять заведомо чрезмерный объём работы.

Для обеспечения отзывчивости интерфейса все потенциально длительные операции, включая серии бенчмарков и криптографические измерения на больших объёмах, выполняются в фоновых горутинах. UI получает обновления по прогрессу и итоговым значениям, что позволяет одновременно наблюдать систему и контролировать ход измерения. Такое разделение особенно важно в режимах, где тест выполняет десятки итераций, а мониторинг параллельно обновляет виджеты с фиксированным периодом. Без подобной конкурентной организации приложение быстро стало бы непрактичным для реального использования.

\subsection{Кросс-платформенные особенности}

В ходе реализации учитываются различия платформ и доступность метрик на разных операционных системах. Хотя библиотека \texttt{gopsutil} и абстрагирует значительную часть платформенно-зависимых деталей, отдельные показатели всё же могут быть недоступны, интерпретироваться по-разному или предоставляться с различной степенью детализации. Это касается, например, средней нагрузки, некоторых счётчиков ввода-вывода, информации о процессах и специфичных характеристик аппаратной платформы.

В таких ситуациях комплекс придерживается принципа деградации функциональности. Если определённая метрика не может быть получена корректно, пользователь продолжает получать остальную доступную информацию, а сама система не прекращает работу и не нарушает выполнение измерительных сценариев. Такой подход особенно важен для прикладного инструмента, который должен быть полезен в разных средах, даже если они предоставляют неидентичный объём телеметрии.

Отдельного упоминания заслуживает поддержка гетерогенных архитектур, в частности платформ семейства Apple Silicon. Для них в проекте предпринята попытка различать производительные и энергоэффективные ядра через системные вызовы \texttt{sysctl}. Даже если такая информация не всегда доступна в полном объёме, сама поддержка подобного сценария подчёркивает ориентацию комплекса на актуальные архитектурные особенности современных вычислительных систем.

\section{Интерпретация результатов и практические рекомендации}

Результаты синтетических тестов в комплексе формируются в прикладных величинах, которые проще соотнести с реальными задачами. Именно этот принцип интерпретации был одним из ключевых при разработке всей системы. Пользователь получает не абстрактный балл, а значение в естественной единице измерения, из которой можно сделать вывод о практической пригодности платформы. Например, пропускная способность AES в МБ/с напрямую отражает потенциальную скорость шифрования трафика в VPN, системе защищённого хранения данных или сервере приложений. Показатель RSA-операций в секунду соотносится с числом возможных асимметричных операций, характерных для аутентификации и рукопожатий TLS. Для хеширования показатель МБ/с показывает, насколько быстро система способна проверять целостность больших архивов, образов, резервных копий или потоков данных.

Схожим образом интерпретируются и другие категории. Процессорные тесты позволяют судить о способности платформы эффективно справляться с вычислительными нагрузками, чувствительными к ветвлениям, рекурсии и масштабированию по потокам. Тесты памяти отражают качество работы с большими линейными массивами и нерегулярными структурами данных. Математические бенчмарки показывают общую эффективность выполнения численных операций и могут быть соотнесены с задачами моделирования, аналитики и научных вычислений.

Для процессорных тестов важна интерпретация не только абсолютного значения, но и стабильности результата. Если коэффициент вариации высок, это означает, что в системе присутствует переменная нагрузка или нестабильный режим исполнения, вмешивающийся в измерения. В таком случае для корректного сравнения платформ рекомендуется:
\begin{itemize}
    \item закрыть ресурсоёмкие приложения и повторить серию измерений;
    \item сопоставить результаты с «топ-процессами» и показателями дискового/сетевого ввода-вывода;
    \item оценить нагрузку по ядрам: равномерная высокая загрузка подтверждает CPU-bound режим, а неравномерная может означать ограничения по синхронизации или приоритетам планировщика.
\end{itemize}

Для подсистемы памяти следует различать ситуации «высокая пропускная способность при высокой латентности» и «низкая пропускная способность при низкой латентности», поскольку в реальных приложениях это приводит к разному поведению. Потоковая обработка данных, мультимедиа и часть научных вычислений чаще зависят от пропускной способности. Базы данных, компиляторы, графовые структуры и индексные хранилища значительно чувствительнее к латентности случайных обращений. Поэтому корректная интерпретация результатов памяти невозможна без понимания прикладного профиля исследуемой задачи.

Показатели дискового и сетевого ввода-вывода также необходимо рассматривать совместно с фоном системы. Высокая скорость чтения или записи не всегда означает устойчиво высокую производительность накопителя: она может быть частично обусловлена системным кэшированием или отсутствием конкурирующих задач в момент измерения. Аналогично, сетевые показатели следует соотносить с внешними условиями и характером текущей активности. Именно поэтому наличие мониторинговой вкладки и списка активных процессов играет не вспомогательную, а методически значимую роль.

В практическом отношении комплекс может использоваться в нескольких сценариях. Он полезен для первичной диагностики рабочей станции или сервера, когда необходимо быстро понять, какой компонент системы ограничивает производительность. Он пригоден для сравнительного анализа аппаратных платформ в учебной и инженерной практике. Наконец, он представляет интерес как демонстрационный инструмент, позволяющий наглядно показать связь между архитектурой вычислительной системы, текущим состоянием ОС и наблюдаемыми результатами синтетических измерений.

Таким образом, реализованный программный комплекс позволяет не только получить числовые показатели, но и связать их с состоянием системы в момент измерения, с качеством статистической устойчивости и с прикладным смыслом результата. Это делает инструмент пригодным как для инженерной диагностики, так и для учебного анализа архитектурных особенностей современных вычислительных платформ и методов программного измерения производительности.
